\documentclass[10pt,twocolumn,a4paper]{article}

\usepackage{xetexko}
\usepackage{geometry}
\usepackage{titlesec}
\usepackage{booktabs}
\usepackage{array}
\usepackage{tabularx}
\usepackage{setspace}
\usepackage{indentfirst}
\usepackage{hyperref}
\usepackage{caption}
\usepackage{float}

\geometry{
  a4paper,
  top=25mm, bottom=25mm,
  left=18mm, right=18mm,
  columnsep=6mm
}

% 한국어 폰트 설정 (macOS 시스템 폰트)
\setmainfont{Apple SD Gothic Neo}[
  BoldFont = Apple SD Gothic Neo Bold,
  ItalicFont = Apple SD Gothic Neo Regular
]
\setsansfont{Apple SD Gothic Neo}
\setmonofont{Courier New}
\setmainhangulfont{Apple SD Gothic Neo}[
  BoldFont = Apple SD Gothic Neo Bold
]

% 제목 서식
\titleformat{\section}[block]
  {\normalsize\bfseries\centering}
  {\Roman{section}.\,}{0.3em}{}
\titlespacing{\section}{0pt}{8pt}{4pt}

\titleformat{\subsection}[block]
  {\normalsize\bfseries}
  {\thesubsection\,}{0.3em}{}
\titlespacing{\subsection}{0pt}{6pt}{3pt}

% 단락 들여쓰기
\setlength{\parindent}{1em}
\setlength{\parskip}{2pt}

% 각주 구분선 제거
\renewcommand{\footnoterule}{\vspace{-3pt}\noindent\rule{0.4\columnwidth}{0.4pt}\vspace{3pt}}

\pagestyle{empty}

\begin{document}

% ── 제목 영역 (전체 폭) ──────────────────────────────────
\twocolumn[{%
\begin{center}
  {\Large\bfseries 실제 관광 경로 기반 AI 에이전트를 활용한}\\[4pt]
  {\Large\bfseries 제주 설화·민담 여행 코스 추천 서비스 개발}\\[10pt]

  {\normalsize 조익준$^\dagger$}\\[2pt]
  {\normalsize 경희대학교 컴퓨터공학과}\\[2pt]
  {\small\itshape e-mail: \{choikjun\}@khu.ac.kr}\\[10pt]

  {\normalsize\bfseries
  Jeju Folklore \& Folktale Travel Course Recommendation Service Development}\\
  {\normalsize\bfseries
  using Real Tourist Route Data and AI Agent}\\[4pt]
  {\normalsize Igjun Choi$^\dagger$}\\
  {\normalsize School of Computing, Kyung Hee University}\\[14pt]
\end{center}
\vspace{2pt}
\noindent\rule{\textwidth}{0.4pt}
\vspace{6pt}
}]

% ── 요약 ────────────────────────────────────────────────
{\small\noindent\textbf{요 약}\quad
본 연구는 제주도 방문 관광객을 대상으로, 실제 관광객이 검증한 여행 코스에
AI가 설화·민담 콘텐츠를 매핑하여 문화 탐험 경험을 제공하는 Flutter
크로스플랫폼 모바일 서비스를 제안한다. 서비스의 중심은 LangGraph StateGraph
기반 코스 추천 에이전트로, 제주관광공사 Visit Jeju 공공 데이터에서 수집한
9,134개의 실제 여행 일정과 6,002개의 관광 장소(GPS 좌표 포함)를 기반으로,
사용자 조건(지역·일수)에 맞는 코스를 선별한 뒤 RAG 검색으로 근처
설화·민담을 매핑하여 문화 콘텐츠가 풍부한 여행 코스를 추천한다.
제주도청 공공 API를 통해 수집한 설화 182건과 민담 323건을 벡터
데이터베이스(ChromaDB)에 적재하며, GPS 좌표와 연결하여 제주 11개 권역에
걸쳐 365건의 위치 기반 콘텐츠를 확보한다. ReAct 패턴의 AI 챗봇과의 대화로
설화를 더 깊이 탐구할 수 있으며, 드라이빙 모드에서는 TTS 해설을 자동
재생한다.}

\vspace{6pt}
{\small\noindent\textbf{키워드}\quad
LangGraph, RAG, 위치 기반 서비스, 제주 설화·민담, Flutter, AI 에이전트}

\vspace{8pt}

% ── I. 서론 ─────────────────────────────────────────────
\section{서 론}

제주도는 연간 1,500만 명 이상이 방문하는 국내 최대 관광지이나, 현행 관광
앱의 대부분은 맛집·숙박·교통 중심의 정보 제공에 머물고 있다. 제주도에는
설화 182건, 민담 323건 등 505건의 문화유산 콘텐츠가 공공 데이터로 개방되어
있음에도, 이를 관광 경험과 결합한 서비스는 부재한 상황이다.

최근 LLM(Large Language Model)의 발전과 함께, 특정 도메인 지식을 벡터
데이터베이스에 저장하고 검색·생성을 결합하는 RAG 기술이 주목받고 있다.
RAG는 모델의 환각(Hallucination) 문제를 억제하면서 도메인 특화 지식을
활용할 수 있어, 정확성이 요구되는 문화유산 콘텐츠 분야에 적합하다.

본 연구는 \textbf{``실제 검증된 여행 코스 위에 설화를 얹는다''} 는 컨셉으로,
제주관광공사 Visit Jeju의 실제 관광객 여행 일정 데이터를 기반으로 LangGraph
AI 에이전트가 설화·민담이 풍부한 코스를 추천하는 Flutter 앱을 구현하는 것을
목표로 한다. 코스의 신뢰성은 실제 관광객 데이터가 보장하고, AI 에이전트는
그 위에 문화 콘텐츠를 매핑·큐레이션하는 역할을 담당한다.

% ── II. 관련 연구 ───────────────────────────────────────
\section{관련 연구}

\subsection{RAG(Retrieval-Augmented Generation)}

RAG는 Lewis et al.(2020)이 제안한 방법으로, 외부 지식 저장소에서 관련
문서를 검색(Retrieve)한 뒤 이를 문맥으로 제공하여 LLM이 더 정확한 답변을
생성(Generate)하도록 한다. 벡터 임베딩 기반의 의미론적 검색을 통해 키워드
검색의 한계를 극복하며, 특정 도메인에 특화된 지식베이스 구축에 활용된다.

\subsection{위치 기반 문화 콘텐츠 서비스}

Pok\'{e}mon GO, Ingress 등의 위치 기반 서비스는 GPS를 활용하여 사용자가
실제 장소를 방문하도록 유도한 바 있다. 문화유산 분야에서는 Google Arts \&
Culture, 국내 문화재청의 e뮤지엄 등이 위치 기반 콘텐츠를 제공하나, 지역
설화·민담과 같은 무형 문화유산과 AI 코스 추천을 결합한 서비스는 미흡한
실정이다.

\subsection{대화형 AI 인터페이스}

ChatGPT, Claude 등의 등장 이후 사용자가 자연어로 정보를 탐색하는 대화형
인터페이스가 보편화되고 있다. RAG를 결합한 도메인 특화 챗봇은 일반 LLM의
지식 한계를 보완하면서 특정 주제에 깊이 있는 답변을 제공한다. 본 연구에서는
이를 관광 맥락에 적용하여 현장형 AI 문화 가이드를 구현한다.

\subsection{LLM 에이전트 및 LangGraph}

Yao et al.(2023)이 제안한 ReAct(Reasoning + Acting) 패턴은 LLM이 추론
단계와 도구 사용 단계를 반복하며 복잡한 작업을 수행하도록 한다. LangGraph는
이러한 에이전트 워크플로우를 유향 그래프(StateGraph)로 명시적으로 정의하여,
조건부 분기·루프·병렬 실행을 지원한다. 단순 체인(Chain) 방식과 달리 중간
결과를 스스로 평가하고 재시도할 수 있어, 다단계 추론이 필요한 코스
선별·검증 태스크에 적합하다.

% ── III. 시스템 설계 ────────────────────────────────────
\section{시스템 설계}

\subsection{전체 아키텍처}

본 시스템은 Flutter 앱, FastAPI 백엔드, 데이터 레이어의 3계층으로 구성된다.
Flutter 앱은 REST/SSE를 통해 백엔드와 통신하며, 백엔드는 LangGraph 코스
추천 에이전트, ReAct 챗봇 에이전트, GPS 핀 조회 API, TTS API로 구성된다.
데이터 레이어는 Visit Jeju 여행일정 DB(9,134개 코스, 6,002개 장소),
ChromaDB(설화·민담 벡터 DB), SQLite(메타데이터·GPS 좌표)로 이루어진다.

사용자 흐름은 계획·경험·기록의 단일 여정으로 구성된다. 계획 단계에서는
지도에서 설화 마커를 탐색하고 지역·일수를 입력하면 LangGraph 에이전트가
코스를 추천한다. 미리보기 화면에서 장소 목록과 매핑된 설화를 확인한 뒤
저장하고, 여행 당일 탐험 모드를 시작하면 GPS 안내와 함께 장소 도착 시
설화가 제공된다. 여행 후에는 멀티모달 스토리를 생성할 수 있다(스트레치
목표).

\subsection{데이터 구조 및 전처리}

\textbf{Visit Jeju 여행 데이터.}\quad 제주관광공사 Visit Jeju 공공데이터에서
두 종류의 데이터를 확보하였다. 여행장소 데이터는 6,002개 관광 장소의
장소명·도로명주소·위도·경도를 포함하며, 여행세부일정 데이터는 146,508행,
9,134개 여행 일정의 일정 ID·장소명·여행 일차·시작 시간·이동 수단을 포함한다.
두 파일을 장소명 기준으로 조인하면 일정별 방문 장소의 GPS 좌표를 확보할 수
있다. 전처리 후 \texttt{courses} 테이블과 \texttt{course\_places} 테이블로
SQLite에 저장한다.

\textbf{설화·민담 GPS 좌표 연결.}\quad 설화(182건)는 텍스트 내 지명을 제주
행정구역 화이트리스트(120개 항목) 기반 매칭으로 추출한 뒤 Geocoding한다.
예비 실험 결과 182건 중 138건(75.8\%)에서 유효 지명이 추출되었으며, 유니크
장소 80곳을 확인하였다. 민담(323건)은 채록 메타데이터의 조사 장소를 파싱하여
Geocoding한다. 통합 결과 제주 11개 권역에 걸쳐 365건의 설화·민담이 GPS와
연결된다.

\subsection{LangGraph 에이전트 기반 AI 코스 추천}

코스 추천은 LangGraph StateGraph로 구현한다. AI가 코스를 처음부터 생성하는
대신, 실제 관광객이 검증한 9,134개의 Visit Jeju 일정 중 조건에 맞는 코스를
선별하고, RAG로 각 코스 장소 인근의 설화·민담을 매핑한다.

StateGraph의 상태는 \{user\_input, region, duration\_days,
candidate\_courses, folklore\_mapped\_courses, ranked\_courses, validated,
retry\_count\}로 정의된다. 노드 구성은 다음과 같다. (1)
\textbf{analyze\_request}: 지역·여행 일수·이동 수단 파악, (2)
\textbf{search\_courses}: Visit Jeju DB에서 조건에 맞는 후보 코스 검색, (3)
\textbf{map\_folklore}: 각 후보 코스의 장소별 반경 내 설화·민담 RAG 매핑,
(4) \textbf{rank\_courses}: 설화 커버리지·다양성 기준 코스 랭킹, (5)
\textbf{validate\_course}: 총 거리·시간·설화 유형 다양성 검증. 검증 통과 시
코스 카드를 생성하며, 실패 시 최대 3회까지 조건을 완화하여 재검색한다.

생성된 코스는 기기 로컬 SQLite에 저장하며, 미리보기 화면에서 지도 위 경로·
방문 장소·매핑된 설화 목록·예상 소요 시간을 확인한 뒤 저장한다.

\subsection{ReAct 에이전트 기반 AI 챗봇 및 드라이빙 오디오}

챗봇은 LangChain ReAct 에이전트로 구현한다. 에이전트는
\texttt{search\_folklore}(설화 RAG 검색), \texttt{search\_folktale}(민담 RAG
검색), \texttt{get\_nearby\_pins}(GPS 핀 조회) 툴을 보유하며, 설화와 민담은
ChromaDB 컬렉션이 분리되어 있어 타입별 검색이 가능하다. max\_distance 임계값
(코사인 거리 0.62)으로 저관련 청크를 제거하여 환각을 억제하며, 출처는 시스템
후처리로 부착한다.

설화·민담 상세 화면에는 TTS 재생 기능을 제공한다. RAG가 생성한 현대 한국어
해설을 OpenAI TTS API로 음성 변환하여 청취할 수 있다. 드라이빙 모드에서는
알림 액션 버튼에서 직접 TTS를 재생할 수 있으며, GPS 핀 도착 시 해당 설화가
자동 재생된다.

\subsection{멀티모달 여행 스토리 생성 (스트레치 목표)}

여행 중 방문한 GPS 장소 이력과 열람한 설화·민담 목록을 수집하여 GPT-4o
프롬프트에 구조화된 여행 문맥으로 제공한다. RAG로 설화 원문을 문맥에
주입하므로 LLM이 내용을 지어내지 않는다. 생성된 에세이의 주요 장면은 DALL-E
API로 삽화를 생성하고, TTS API로 음성 변환하여 팟캐스트 형태로 제공한다.

% ── IV. 학기 진행 계획 ──────────────────────────────────
\section{학기 진행 계획}

\begin{table}[H]
\centering
\caption{학기 구현 일정}
\small
\begin{tabularx}{\columnwidth}{@{}lX@{}}
\toprule
\textbf{기간} & \textbf{주요 내용} \\
\midrule
1\textasciitilde2주차 & Visit Jeju 데이터 전처리, 장소명 조인 및 courses DB 구축, 설화·민담 GPS 좌표 완성 \\
3\textasciitilde4주차 & FastAPI 백엔드 구조, LangGraph StateGraph 코스 추천 에이전트 구현 \\
5\textasciitilde7주차 & Flutter 앱: 홈 지도 화면(마커·클러스터링), 코스 추천·미리보기·저장 화면 \\
8\textasciitilde10주차 & Flutter 앱: 탐험 모드(GPS, 푸시 알림, 설화 상세, TTS, 드라이빙 모드) \\
11\textasciitilde12주차 & ReAct 챗봇 에이전트 구현, Flutter 연동 및 질의응답 인터페이스 \\
13\textasciitilde14주차 & (스트레치) 멀티모달 스토리 생성 구현 및 통합 테스트 \\
15주차 & 사용자 평가 실험 및 결과 분석 \\
16주차 & 최종 보고서 작성 및 발표 준비 \\
\bottomrule
\end{tabularx}
\end{table}

% ── V. 결론 ─────────────────────────────────────────────
\section{기대 효과 및 결론}

본 연구는 실제 관광객 여행 데이터와 RAG·LangGraph 에이전트를 결합하여
제주 설화·민담이라는 무형 문화유산의 접근성을 높이는 것을 목표로 한다.
AI가 코스를 처음부터 창작하는 대신, Visit Jeju의 9,134개 실제 검증 코스를
기반으로 설화 매핑에 집중함으로써 코스의 현실성과 AI 큐레이션의 가치를
동시에 확보한다. 코스 추천(계획)·GPS 탐험(경험)·스토리 기록(기억)이 하나의
여정으로 연결되며, 설화는 텍스트 지명 추출, 민담은 채록 장소 활용이라는
이중 전략으로 505건 전체를 GPS 기반 서비스에 통합한다.

향후 연구로는 다국어 지원을 통한 외국인 관광객 대상 확장, 사용자 행동 기반
개인화 추천을 고려한다.

% ── 참고문헌 ────────────────────────────────────────────
\section*{참 고 문 헌}

\begin{small}
\begin{enumerate}
\setlength{\itemsep}{1pt}
\item P. Lewis et al., ``Retrieval-Augmented Generation for
  Knowledge-Intensive NLP Tasks,'' \textit{NeurIPS}, 2020.
\item 제주특별자치도, 제주 설화정보 공공데이터 API (E07), 2023.
\item 제주특별자치도, 제주 민담정보 공공데이터 API (E08), 2023.
\item 제주관광공사, 제주관광정보시스템(VISIT JEJU) 여행장소 및 여행세부일정
  공공데이터, 2026.
\item Chroma, ``Chroma: the AI-native open-source embedding database,''
  2023.
\item OpenAI, ``Text Embedding Models,''
  \textit{platform.openai.com/docs/guides/embeddings}, 2024.
\item Google, ``Google Maps Platform -- Geocoding API \& Directions
  API,'' \textit{developers.google.com/maps}, 2024.
\item S. Yao et al., ``ReAct: Synergizing Reasoning and Acting in
  Language Models,'' \textit{ICLR}, 2023.
\item LangChain, ``LangGraph: Build Stateful Multi-Actor
  Applications,'' \textit{langchain-ai.github.io/langgraph}, 2024.
\end{enumerate}
\end{small}

\vspace{6pt}
\noindent\rule{\columnwidth}{0.4pt}
\vspace{4pt}
{\footnotesize $^\dagger$ 본 연구는 경희대학교 컴퓨터공학과 졸업프로젝트로 수행되었음.}

\end{document}
